\documentclass[10pt,a4paper]{amsart}
\usepackage[margin=19mm]{geometry}
\usepackage{amsmath,amssymb,mathtools}
\usepackage[scr=rsfs,cal=euler]{mathalfa}
\usepackage{xfrac}
\usepackage[unicode,hidelinks,bookmarks=false]{hyperref}
\newcommand{\ie}{i.e.\ }
\newcommand{\cf}{cf.\ }
\newcommand{\resp}{respectively\ }
\newcommand{\Subsection}{\subsection}
\providecommand{\nobreakdash}{\nobreak\hbox{-}}
\setlength{\emergencystretch}{2em}
\multlinegap0pt
\usepackage{tikz}
\usepackage{caption}
\usepackage{float}
\newtheorem{lemma}{Lemma}
\newtheorem{theorem}{Theorem}
\title{Signs of Hamiltonian Circles in Simple Plane Signed Graphs}
\author{Xiyong Yan}
\date{Source: arXiv:2602.21530v1 (CC BY 4.0)}
\begin{document}
\maketitle

\noindent\emph{Source and licence.} Signs of Hamiltonian Circles in Simple Plane Signed Graphs. Version arXiv:2602.21530v1 (CC BY 4.0), distributed by arXiv under the Creative Commons Attribution 4.0 International licence, http://creativecommons.org/licenses/by/4.0/, as recorded in the arXiv licence metadata for this exact version and shown on the abstract page https://arxiv.org/abs/2602.21530v1. Full licence notice: \texttt{licenses/arxiv-2602.21530-CC-BY-4.0.txt}.\par\smallskip

\noindent\textit{Editorial scope.} Counted unit: the article's theorem thm:all-same-sign (source0.tex lines 984--987). The article's complete original proof of it is delivered with it (source0.tex lines 990--1124, one \texttt{proof} environment of 135 lines). No mathematics is written, repaired or re-derived: the statement and the proof are the article's own lines, in the article's own notation, and no retained line is substituted or re-flowed.  Retained uncounted as the source-local context the proof needs: lines 527--540; lines 937--948.  Nothing was omitted from any retained span, so the package carries no omission record.  No retained line contains a citation, so the wrapper carries no bibliography and no bibliography entry.  No retained line contains a figure environment, an image-inclusion command or drawing code, so no image is delivered and the package declares no asset dependency.  Editorial formatting only, in addition: an AMS \texttt{amsart} preamble that restates the article's own declarations and macros used by the retained lines, namely the environments \texttt{lemma}, \texttt{theorem} on separate counters, together with the standard packages those lines need.  The retained environments print in this document as Lemma 1, Theorem 1 in order of appearance, whereas the article prints the same items with the numbering of its own full document.  The complete native source of the article as distributed in the arXiv e-print for this exact version is bundled unchanged as \texttt{sources/195272/source0.tex}.
\begin{lemma}\label{facegraph}
Let $\Sigma$ be a $2$-connected simple plane signed graph, let $H'$ be the set of all bounded faces of $\Sigma$,
and let $C_0$ be the boundary circle of the outer face.
Let $\sigma:E(G)\to\{+,-\}$ be an edge-signing.

For each bounded face $f\in H'$, define its face-sign by
\[
\sigma(f)=\prod_{g\in \partial f}\sigma(g).
\]
Then
\[
\sigma(C_0)=\prod_{f\in H'} \sigma(f).
\]
\end{lemma}

\begin{lemma}\label{lemma123} 
Let $\Sigma$ be an $m\times n$ grid with $m,n>3$ and $n$ even.
Let $H_1$ and $H_2$ be two distinct Hamiltonian circles with corresponding Hamiltonian sets
$H_1'$ and $H_2'$, respectively. Suppose
\[
H_1' \triangle H_2'=\{a,b\},
\]
where $a\in H_1'\setminus H_2'$ and $b\in H_2'\setminus H_1'$ are boxes. Then
\[
\sigma(H_1)=\sigma(H_2)\quad \Longleftrightarrow\quad \sigma(a)=\sigma(b).
\]
\end{lemma}

\begin{theorem}\label{thm:all-same-sign}
Let $\Sigma$ be an $m\times n$ grid with $m$ even and $m,n>3$.
 All Hamiltonian circles in $\Sigma$ have the same sign if and only if all boxes of $\Sigma$ except the four corner boxes have the same sign.
\end{theorem}

\begin{proof}
Suppose all Hamiltonian circles in $\Sigma$  have the same sign.



We will show  that all boxes of $\Sigma$, except the four corner boxes, have the same sign.
We prove this direction in two cases.

\medskip
\noindent\emph{Case 1: $n$ is even.}
Let $\mathcal{B}$ denote the set of all bounded faces of the graph $\Sigma$.

Let
\[
L:=\{[i,j]\mid i\in \{1,2,\dots,m-2\},\ j\in \{2,4,\dots,n-2\}\}
\]
be a co-Hamiltonian set of $\mathcal{B}$, and let $H=\mathcal{B}\setminus L$.
Fix a column index  $j\in\{2,4,\dots,n-2\}$, and each $i\in\{1,2,\dots,m-2\}$,
we may replace the box $[m-1,j]$ in $H$ by the box $[i,j]$
(that is, delete $[m-1,j]$ and add $[i,j]$).
Denote the resulting Hamiltonian set by $H_{i,j}$.
 Since the Hamiltonian circle in $H_{i,j}$
and the Hamiltonian circle in $H$ have the same signs and $H_{i,j}\triangle H=\{[i,j],[m-1,j]\}$ , by Lemma \ref{lemma123}, $\sigma([i,j])=\sigma([m-1,j])$. Since $i$ is arbitrary, all the boxes in the column $j$ have the same sign. 


Consider $H$ again. Delete the box $[2,1]$ and add the box $[1,2]$
(equivalently, move the box $[2,1]$ to position $[1,2]$).
Denote the resulting Hamiltonian set by $H_2$.
Then, in $H_2$, for any $k\in\{3,4,\dots,m-2\}$, we may replace the box $[k,1]$
by the box $[2,1]$ (that is, move $[k,1]$ to position $[2,1]$).
Denote the resulting Hamiltonian set by $H_{k,1}$. Since all the Hamiltonian circles have the same sign and $H_{k,1}\triangle H_2=\{[2,1],[k,1]\}$, by Lemma \ref{lemma123}, $\sigma([k,1])=\sigma([2,1])$. Hence, all non-corner boxes in the first column must have the same sign. Similarly, all non-corner boxes in the last column must have the same sign.



Next, consider $\Sigma$ again. Let
\[
L_2:= \{[i,j]\mid i\in \{1,2,\dots,m-2\},\ j\in \{3,5,\dots,n-3\}\}
\;\cup\;
\{[i,1],[i,n-1]\mid i\in \{2,4,\dots,m-2\}\}.
\]
Let $H'=\mathcal{B}\setminus L_2$.
Then $H'$ contains exactly one Hamiltonian circle.

\medskip
\noindent\emph{Row movements in $H'$.}
Fix an odd column index $j\in\{3,5,\dots,n-3\}$.
For any $i\in\{1,2,\dots,m-2\}$, replace the box $[m-1,j]$ in $H'$
by the box $[i,j]$, and denote the resulting Hamiltonian set by $H_{i,j}$. Since we assume that all Hamiltonian circles in $\Sigma$ have the same sign, by Lemma~\ref{lemma123} we have
$\sigma([i,j])=\sigma([m-1,j])$.
Since $i$ is arbitrary, we conclude that all boxes in column $j$ have the same sign.


\medskip
\noindent\emph{Relations between different columns.}
Finally, we compare boxes lying in different columns.
In $H$, fix a number $j\in\{1,3,\dots,n-3\}$, we may move the box $[2,j]$
to position $[1,j+1]$. Denote the resulting Hamiltonian set by $H_{1,j+1}$. Since all the Hamiltonian circles have the same sign, 
by Lemma \ref{lemma123}, $\sigma([2,j])=\sigma([1,j+1])$. This implies the non-corner box sign in column $l$ is the same as the non-corner box sign in column $l+1$, for $l$ is odd.

Similarly, in $H$ again,  for each $j\in\{3,5,\dots,n-1\}$, we may move the box $[2,j]$
to position $[1,j-1]$. Denote the resulting Hamiltonian set by $H_{1,j-1}$.
With the same reasoning, $\sigma([2,j])=\sigma([1,j-1])$.  This implies the non-corner box sign in column $l$ is the same as the non-corner box sign in column $l+1$, for $l$ is even.
Hence, any two non-corner box in different column will have the same sign. 




\medskip
\noindent\emph{Case 2: $n$ is odd.}

Let
\[
L:=\{[i,j]\mid i\in \{1,2,\dots,m-2\},\ j\in \{2,4,\dots,n-3\}\}
\;\cup\;
\{[i,n-1]\mid i\in \{2,4,\dots,m-2\}\}
\]
be a co-Hamiltonian set of $\Sigma$, and let $H=\mathcal{B}\setminus L$.
Then $H$ contains exactly one Hamiltonian circle.

Using the same row and column movement arguments as in Case~1,
we obtain that all non-corner boxes in the first $n-2$ columns
must have the same sign since all Hamiltonian circles in $\Sigma$
to have the same sign.

Next, consider a different co-Hamiltonian set
\[
L':=\{[i,j]\mid i\in \{1,2,\dots,m-2\},\ j\in \{3,5,\dots,n-2\}\}
\;\cup\;
\{[i,1]\mid i\in \{2,4,\dots,m-2\}\}.
\]
Using the same row and column movement arguments as in Case~1,
we conclude that all non-corner boxes in columns $2,3,\dots,n-1$
must have the same sign, since all Hamiltonian circles in $\Sigma$
are assumed to have the same sign.

Hence, from both cases we conclude that all boxes of $\Sigma$, except the four corner boxes, must have the same sign.




Conversely, assume that all boxes of $\Sigma$, except the four corner boxes, have the same sign.
We show that all Hamiltonian circles in $\Sigma$ have the same sign.

From the forward direction, we make the following two observations:
\begin{enumerate}
\item every Hamiltonian set contains all four corner boxes;
\item all Hamiltonian sets contain the same number of boxes.
\end{enumerate}

Let $H'$ be any Hamiltonian set, and let $C$ denote its corresponding Hamiltonian circle.
By Lemma \ref{facegraph},
\[
\sigma(C)=\prod_{[i,j]\in H'} \sigma([i,j]).
\]
Write
\[
H' = \{\text{four corner boxes}\} \cup \{\text{non-corner boxes in } H'\}.
\]

Since all non-corner boxes have the same sign, say $\alpha$, and since every Hamiltonian set
contains the same number of non-corner boxes, the contribution of the non-corner boxes to
$\sigma(C)$ is the same for every Hamiltonian circle.
Moreover, the contribution of the four corner boxes is fixed, since every Hamiltonian set
contains all four of them.

Therefore, for any two Hamiltonian circles $C_1$ and $C_2$, we have
\[
\sigma(C_1)=\sigma(C_2).
\]
Hence, all Hamiltonian circles in $\Sigma$ have the same sign.




\end{proof}

\end{document}
